% Part: counterfactuals
% Chapter: minimal-change-semantics
% Section: transitivity

\documentclass[../../../include/open-logic-section]{subfiles}

\begin{document}

\olfileid{cnt}{min}{cpo}

\olsection{عکسِ نقیض}

مادی تے سخت شرطیے اپنے عکسِ نقیض دے ہم ارز نیں۔ خلافِ واقعہ شرطیاں لئی
ایہ گل سچ نہیں۔ کریٹزر دی اک مثال ویکھو:
\begin{quote}
  جے گوئٹے 1832 وچ نہ مویا ہوندا، تاں ہُن (فیر وی) مریا ہویا ہوندا۔

  جے گوئٹے ہُن مریا ہویا نہ ہوندا، تاں اوہ 1832 وچ مویا ہوندا۔
\end{quote}
پہلا جملہ سچ اے، کیوں جے انسان سیکڑیاں ورھے نہیں جیندے۔ دوجا صاف جھوٹ اے۔
جے گوئٹے ہُن مریا ہویا نہ ہوندا، تاں اوہ ہجے جیوندا ہوندا، تے ایس لئی
1832 وچ نہیں مر سکیا ہوندا۔

\begin{ex}\ollabel{ex:contraposition-counterex}
  گولیاں والی معنیات وچ عکسِ نقیض غیر معتبر اے، یعنی ساڈے کول $p
  \cif q \Entails/ \lnot q \cif \lnot p$ اے۔ $p$ دا مطلب ’گوئٹے
  1832 وچ نہیں مویا‘ تے $q$ دا مطلب ’گوئٹے ہُن مریا ہویا اے‘ لو۔
  ایہ گل نمونے $\mModel{M_1} = \tuple{W, O, V}$ وچ دکھا سکدے آں، جیہدے وچ
  $W = \{w, w_1, w_2\}$، $O = \{\{w\}, \{w, w_1\}, \{w, w_1, w_2\}\}$،
  $V(p) = \{w_1, w_2\}$ تے $V(q) = \{w, w_1\}$۔ سو $w$ اصل جہان اے
  جتھے گوئٹے 1832 وچ مویا تے ہجے وی مریا ہویا اے۔ $w_1$ اوہ (نیڑلا)
  جہان اے جتھے گوئٹے، منّ لو، 1833 وچ مویا تے ہجے وی مریا ہویا اے۔
  تے $w_2$ اوہ (دُور دا) جہان اے جتھے گوئٹے ہجے وی جیوندا اے۔
\begin{figure}
\begin{center}
\begin{tikzpicture}[scale=.6,layerwidth=1.5]\tiny
  \spheresystem{3}
  \begin{scope}
    \clip \propositionplot{0}{.8}{50}{4.5} -- (4.5,-4.5) -- (-4.5,-4.5) -- (-4.5,4.5) -- (4.5,4.5);
    \spherelayer{3}
  \end{scope}
  \propositionintersect{0}{2}{110}{5.5}
  \proposition{0}{.8}{50}{4.5}
  \path (0,0) node[world] {$w$};
  \spherepos{0}{2}{node[world] {$w_1$}}
  \spherepos{-50}{3}{node[world] {$w_2$}}
  \spherepos{28}{3}{node {$q$}}
  \spherepos{42}{3}{node {$\lnot q$}}
  \spherepos{-55}{4}{node {$p$}}
  \spherepos{-67}{4}{node {$\lnot p$}}
\end{tikzpicture}
\caption{عکسِ نقیض دی ردّی مثال}
\ollabel{fig:contraposition}
\end{center}
\end{figure}
$p$ نوں تھاں دین والا گولا $S = \{w, w_1\}$ موجود اے تے $p \lif q$
اوہدے سارے جہاناں اُتے سچ اے، سو $\mSat{M}{p \cif q}[w]$۔ پر $\lnot q$
نوں تھاں دین والا گولا $\{w, w_1, w_2\}$ وچ اک ایہو جیہا جہان اے،
یعنی~$w_2$، جتھے $q$ جھوٹ تے $p$ سچ اے، سو $\mSat/{M}{\lnot q \lif \lnot p}[w_2]$۔
\end{ex}

\end{document}
